\cvsection{Expériences professionnelles}

\begin{cventries}
  \cventry
    {Stagiaire ingénieur de recherche en IA}
    {INRIA}
    {Bordeaux, France}
    {Mars 2026 -- Septembre 2026}
    {
      \begin{cvitems}
        \item {Transformation d’un besoin utilisateur en une application fonctionnelle aidant les chercheurs à trouver des articles scientifiques.}
        \item {Automatisation de la collecte de métadonnées d’articles scientifiques en ligne et transmission des résultats de tâches cron par webhooks.}
        \item {Développement du backend Python/FastAPI avec API REST, OpenAPI et MongoDB, puis d’une interface React/Next.js.}
        \item {Mise en œuvre d’un pipeline RAG avec préparation des documents, embeddings, recherche sémantique et stockage vectoriel dans ChromaDB.}
        \item {Gestion d’une application auto-hébergée utilisant des LLM locaux, avec installation, configuration et utilisation de vLLM sur GPU NVIDIA.}
        \item {Comparaison de Qwen, Gemma et GLM et de plusieurs quantifications selon la qualité, la latence, le TTFT et la taille du contexte.}
        \item {Conteneurisation de composants avec Docker et mise en place d’un pipeline CI/CD sur GitHub.}
        \item {Évaluation d’OpenClaw dans un environnement Docker isolé avec une attention portée aux risques liés à l’accès autonome aux outils et aux fichiers.}
        \item {Définition du protocole expérimental, analyse des résultats et contribution à un article scientifique comparant les modèles.}
      \end{cvitems}
    }

  \cventry
    {Stagiaire de recherche en apprentissage par renforcement}
    {Université Karunya}
    {Coimbatore, Tamil Nadu, Inde}
    {Juin 2025 -- Septembre 2025}
    {
      \begin{cvitems}
        \item {Développement d’un modèle d’apprentissage par renforcement et de son environnement d’entraînement pour la navigation autonome d’un robot hospitalier.}
        \item {Déploiement et optimisation du logiciel et du modèle sur un Raspberry Pi embarqué.}
      \end{cvitems}
    }
\end{cventries}
